\documentclass{apex_mdd}
\modelid{APEX-MDL-0007}
\mddversion{1.0}
\mddowner{Dr.\ Priya Nair (CRO)}
\mddvalidator{Dr.\ Samuel Achebe (Model Validation Officer)}
\mddstatus{Validated}
\reviewdate{2026-03-01}
\nextreview{2027-03-01}
\regulatory{IFRS 9.5.5, EBA/GL/2020/06, SR 11-7}

\title{IFRS 9 Expected Credit Loss Model}
\begin{document}
\maketitle
\statusbadge

\tableofcontents
\newpage

\section{Model Overview}

\subsection{Purpose}
The IFRS 9 ECL model computes Expected Credit Loss provisions for the loan portfolio under the three-stage impairment framework. Provisions are recognised on a 12-month basis (Stage 1) or lifetime basis (Stage 2/3) depending on significant increases in credit risk (SICR).

\subsection{Business Use}
\begin{itemize}
  \item \textbf{Financial reporting:} ECL provision recorded on income statement; affects CET1 ratio via retained earnings.
  \item \textbf{Credit risk management:} Stage classification drives credit officer workout strategy.
  \item \textbf{Capital planning:} Stressed ECL feeds DFAST adverse scenario P\&L projections.
  \item \textbf{Pricing:} Risk-adjusted return on equity uses ECL as expected loss charge.
\end{itemize}

\section{Theoretical Basis}

\subsection{ECL Formula}
\begin{formulabox}
\[
  \mathrm{ECL} = PD \times LGD \times EAD \times \mathrm{DF}
\]
For each obligor $i$ across periods $t$:
\[
  \mathrm{ECL}_i = \sum_{t=1}^{T_i} PD_i(t\mid t-1) \times LGD_i \times EAD_i(t) \times \mathrm{DF}(t)
\]
\end{formulabox}

\subsection{Stage Definitions}
\begin{formulabox}
\textbf{Stage 1} (no SICR): $\quad\mathrm{ECL}_{S1} = PD_{12m} \times LGD \times EAD$\\[4pt]
\textbf{Stage 2} (SICR): $\quad\mathrm{ECL}_{S2} = \sum_{t=1}^{T_{\mathrm{life}}} PD_{\mathrm{marg}}(t) \times LGD \times EAD(t) \times \mathrm{DF}(t)$\\[4pt]
\textbf{Stage 3} (default, $>$90 DPD): $\quad\mathrm{ECL}_{S3} = LGD \times EAD_{\mathrm{default}}$
\end{formulabox}

\subsection{SICR Definition}
An obligor is reclassified Stage 1$\to$2 if \textit{any} of:
\begin{itemize}
  \item Credit rating downgrade $\geq 2$ notches since origination.
  \item Days past due $> 30$.
  \item Watch-list or watch-list-equivalent designation.
  \item Qualitative override (credit officer judgement).
\end{itemize}

\subsection{PD Estimation}
Through-the-cycle (TTC) PD from Moody's idealized cumulative default rates, adjusted via macro satellite model:
\begin{formulabox}
\[
  PD_{\mathrm{PIT}}(t) = PD_{\mathrm{TTC}} \times \exp\!\bigl(\beta_{\mathrm{GDP}}\,\Delta\mathrm{GDP}(t) + \beta_{\mathrm{UR}}\,\Delta\mathrm{UR}(t)\bigr)
\]
$\beta_{\mathrm{GDP}} = -0.8$ (GDP growth $+1\% \Rightarrow PD -8\%$),\quad $\beta_{\mathrm{UR}} = +0.6$ (unemployment $+1\mathrm{pp} \Rightarrow PD +6\%$).
\end{formulabox}

\subsection{LGD and EAD}
\textbf{LGD:}
\[
  LGD = 1 - \text{Recovery},\qquad \text{Recovery} = \min\!\left(\frac{\text{Collateral}}{EAD},\;0.90\right)\times\text{haircut}
\]
Unsecured: $LGD = 0.45$ (corporate), $0.25$ (secured mortgages). Haircuts: CRE 30\%, residential 10\%.

\textbf{EAD:} $EAD = \text{drawn} + CCF \times \text{undrawn}$, where $CCF = 0.75$ (revolving), $0.50$ (term).

\section{Mathematical Specification}

\textbf{Sample portfolio:} 50 obligors.
\begin{center}
\small
\begin{tabular}{ll}
\toprule
\textbf{Parameter} & \textbf{Value} \\
\midrule
Rating distribution & 60\% IG (BBB+/BBB/BBB$-$), 30\% BB, 10\% B/CCC \\
Average PD (Stage 1) & 1.2\% \\
Average LGD & 42\% \\
ECL coverage (normal) & $\sim$1.8\% \\
ECL coverage (stressed) & $\sim$3.5\% \\
\bottomrule
\end{tabular}
\end{center}

\section{Implementation}

\textbf{Code location:} \texttt{infrastructure/credit/ifrs9\_ecl.py}\\
\textbf{Class:} \texttt{IFRS9ECLEngine}\\[4pt]
\textbf{Key methods:}
\begin{itemize}
  \item \texttt{calculate\_portfolio\_ecl()} --- full portfolio ECL.
  \item \texttt{classify\_stages()} --- SICR assessment.
  \item \texttt{run\_scenario(macro\_params)} --- stressed ECL for DFAST.
\end{itemize}
\textbf{Runtime:} $\sim$5\,ms for 50 obligors.

\section{Validation}

\textbf{Stage classification:} Validated against sample obligors with known PD migration history; accuracy 94\%.

\textbf{Macro sensitivity:} $\beta_{\mathrm{GDP}}$ validated against EBA stress test results for comparable European banks; within $\pm 15\%$.

\textbf{Coverage ratio:} 1.8\% (normal) vs EBA peer average 1.6\%; slightly conservative, acceptable.

\section{Model Limitations}

\begin{enumerate}
  \item \textbf{TTC PD basis:} TTC PD underestimates current default risk in downturns; PIT adjustment partially corrects this.
  \item \textbf{50-obligor sample:} Production portfolio has 850 obligors; demo portfolio is illustrative only.
  \item \textbf{Flat LGD:} Collateral values not mark-to-market; static LGD may understate loss in collateral price declines.
  \item \textbf{Macro satellite not post-COVID calibrated:} GDP/unemployment coefficients estimated pre-2020.
\end{enumerate}

\section{Open Findings}

\begin{findings}
\finding{ECL-F1}{Major}{Macro satellite model GDP coefficient ($\beta_{\mathrm{GDP}} = -0.8$) not recalibrated post-COVID; EBA/GL/2020/06 requires calibration to reflect current economic uncertainty.}{Open}
\end{findings}

\end{document}
